% GENERATED FILE. Edit canonical lessons, then run npm run generate:books.
% canonical-source-hash: fnv1a64:935bb98cc60c1620
% canonical-lessons: TE-C439-rijarvesanu, TE-C439-bukingu, TE-C439-bertu, TE-C439-bogi, TE-C439-kirayi

\chapter{Booking a Berth}
\label{ch:te-booking-a-berth}
\hlchaptermodality{439}

\begin{chapteropening}
\textbf{By the end of this chapter:} \emph{I can say a reservation, a booking, a berth, a railway coach and a fare.}

Complete the last lesson of chapter 439: I can say a reservation, a booking, a berth, a railway coach and a fare.
\end{chapteropening}

\section[\texorpdfstring{\te{రిజర్వేషను} (rijarvēṣanu)}{రిజర్వేషను (rijarvēṣanu)}]{\te{రిజర్వేషను} (rijarvēṣanu) — a reservation}

\label{lesson:TE-C439-rijarvesanu}



\begin{quote}
Before the new one: say the Telugu for a route, then the Telugu for to set foot in.
\end{quote}

\subsection*{You'll want to know: \te{రిజర్వేషను}}
\textbf{\te{రిజర్వేషను}} — \emph{rijarvēṣanu} — \textquotedblleft{}a reservation\textquotedblright{}.

\begin{quote}
\te{పండుగ} \te{రోజుల్లో} \te{రైలులో} \te{రిజర్వేషను} \te{దొరకడం} \te{కష్టం}.
\end{quote}

\begin{quote}
\emph{paṇḍuga rōjullō railulō rijarvēṣanu dorakaḍaṁ kaṣṭaṁ.}
\end{quote}

\begin{quote}
In the festival season it is hard to get a reservation on the train.
\end{quote}

Like most English loans in Telugu it ends in -\te{ఉ}: \te{రిజర్వేషను}. It goes with \te{చేయు}: \te{రిజర్వేషను} \te{చేయించాను}, I had a seat reserved.

\begin{cousinweb}
\textbf{\te{రిజర్వేషను}} (\emph{rijarvēṣanu}) is a loan from English \emph{reservation}, a seat kept for one passenger. Telugu adds its usual -\te{ఉ} to the final consonant.
\end{cousinweb}

\begin{grammarlens}[title={What you've built}]
One more word for getting around.
\end{grammarlens}

\subsection*{Your turn}
\begin{itemize}
\raggedright
  \item \emph{Say it:} \emph{rijarvēṣanu}
  \item \emph{Say it:} \emph{rijarvēṣanu}, once more
  \item \emph{Say it:} \emph{paṇḍuga rōjullō railulō rijarvēṣanu dorakaḍaṁ kaṣṭaṁ.} — In the festival season it is hard to get a reservation on the train.
  \item \emph{Recall:} say the Telugu for a route, then the Telugu for to set foot in, then say \emph{rijarvēṣanu} again
\end{itemize}

\begin{tcolorbox}[breakable,colback=teal!4,colframe=teal!35!black,title={Before you move on}]
What does \te{రిజర్వేషను} mean? (\textquotedblleft{}A reservation\textquotedblright{}.) Say it once more.
\end{tcolorbox}
\section[\texorpdfstring{\te{బుకింగు} (bukiṅgu)}{బుకింగు (bukiṅgu)}]{\te{బుకింగు} (bukiṅgu) — a booking}

\label{lesson:TE-C439-bukingu}



\begin{quote}
Before the new one: say the Telugu for to set foot in, then the Telugu for a reservation.
\end{quote}

\subsection*{You'll want to know: \te{బుకింగు}}
\textbf{\te{బుకింగు}} — \emph{bukiṅgu} — \textquotedblleft{}a booking\textquotedblright{}.

\begin{quote}
\te{మా} \te{బుకింగు} \te{నిన్ననే} \te{అయిపోయింది}.
\end{quote}

\begin{quote}
\emph{mā bukiṅgu ninnanē ayipōyindi.}
\end{quote}

\begin{quote}
Our booking was done yesterday itself.
\end{quote}

A ticket counter for advance tickets is called the \te{బుకింగు} office in many towns.

\begin{cousinweb}
\textbf{\te{బుకింగు}} (\emph{bukiṅgu}) is a loan from English \emph{booking}, with the Telugu final -\te{ఉ}.
\end{cousinweb}

\begin{grammarlens}[title={What you've built}]
One more word for getting around.
\end{grammarlens}

\subsection*{Your turn}
\begin{itemize}
\raggedright
  \item \emph{Say it:} \emph{bukiṅgu}
  \item \emph{Say it:} \emph{bukiṅgu}, once more
  \item \emph{Say it:} \emph{mā bukiṅgu ninnanē ayipōyindi.} — Our booking was done yesterday itself.
  \item \emph{Recall:} say the Telugu for to set foot in, then the Telugu for a reservation, then say \emph{bukiṅgu} again
\end{itemize}

\begin{tcolorbox}[breakable,colback=teal!4,colframe=teal!35!black,title={Before you move on}]
What does \te{బుకింగు} mean? (\textquotedblleft{}A booking\textquotedblright{}.) Say it once more.
\end{tcolorbox}
\section[\texorpdfstring{\te{బెర్తు} (bertu)}{బెర్తు (bertu)}]{\te{బెర్తు} (bertu) — a sleeper berth}

\label{lesson:TE-C439-bertu}



\begin{quote}
Before the new one: say the Telugu for a reservation, then the Telugu for a booking.
\end{quote}

\subsection*{You'll want to know: \te{బెర్తు}}
\textbf{\te{బెర్తు}} — \emph{bertu} — \textquotedblleft{}a sleeper berth\textquotedblright{}.

\begin{quote}
\te{నాకు} \te{కింది} \te{బెర్తు} \te{కావాలి}.
\end{quote}

\begin{quote}
\emph{nāku kindi bertu kāvāli.}
\end{quote}

\begin{quote}
I would like the lower berth.
\end{quote}

The berths on a sleeper coach are \te{కింది}, the lower, \te{మధ్య}, the middle, and \te{పై}, the upper.

\begin{cousinweb}
\textbf{\te{బెర్తు}} (\emph{bertu}) is a loan from English \emph{berth}, a bed on a train or a ship.
\end{cousinweb}

\begin{grammarlens}[title={What you've built}]
One more word for getting around.
\end{grammarlens}

\subsection*{Your turn}
\begin{itemize}
\raggedright
  \item \emph{Say it:} \emph{bertu}
  \item \emph{Say it:} \emph{bertu}, once more
  \item \emph{Say it:} \emph{nāku kindi bertu kāvāli.} — I would like the lower berth.
  \item \emph{Recall:} say the Telugu for a reservation, then the Telugu for a booking, then say \emph{bertu} again
\end{itemize}

\begin{tcolorbox}[breakable,colback=teal!4,colframe=teal!35!black,title={Before you move on}]
What does \te{బెర్తు} mean? (\textquotedblleft{}A sleeper berth\textquotedblright{}.) Say it once more.
\end{tcolorbox}
\section[\texorpdfstring{\te{బోగీ} (bōgī)}{బోగీ (bōgī)}]{\te{బోగీ} (bōgī) — a railway coach}

\label{lesson:TE-C439-bogi}



\begin{quote}
Before the new one: say the Telugu for a booking, then the Telugu for a sleeper berth.
\end{quote}

\subsection*{You'll want to know: \te{బోగీ}}
\textbf{\te{బోగీ}} — \emph{bōgī} — \textquotedblleft{}a railway coach\textquotedblright{}.

\begin{quote}
\te{మా} \te{బోగీ} \te{రైలు} \te{చివర} \te{ఉంది}.
\end{quote}

\begin{quote}
\emph{mā bōgī railu civara undi.}
\end{quote}

\begin{quote}
Our coach is at the end of the train.
\end{quote}

Plural: \te{బోగీలు}, \emph{bōgīlu}.

The coaches of a train are its \te{బోగీలు}; the number of yours is on the ticket.

\begin{cousinweb}
\textbf{\te{బోగీ}} (\emph{bōgī}) is a loan from English \emph{bogie}, which in Indian English names a railway carriage.
\end{cousinweb}

\begin{grammarlens}[title={What you've built}]
One more word for getting around.
\end{grammarlens}

\subsection*{Your turn}
\begin{itemize}
\raggedright
  \item \emph{Say it:} \emph{bōgī}
  \item \emph{Say it:} \emph{bōgī}, once more
  \item \emph{Say it:} \emph{mā bōgī railu civara undi.} — Our coach is at the end of the train.
  \item \emph{Recall:} say the Telugu for a booking, then the Telugu for a sleeper berth, then say \emph{bōgī} again
\end{itemize}

\begin{tcolorbox}[breakable,colback=teal!4,colframe=teal!35!black,title={Before you move on}]
What does \te{బోగీ} mean? (\textquotedblleft{}A railway coach\textquotedblright{}.) Say it once more.
\end{tcolorbox}
\section[\texorpdfstring{\te{కిరాయి} (kirāyi)}{కిరాయి (kirāyi)}]{\te{కిరాయి} (kirāyi) — a fare; hire charge}

\label{lesson:TE-C439-kirayi}



\begin{quote}
Before the new one: say the Telugu for a sleeper berth, then the Telugu for a railway coach.
\end{quote}

\subsection*{You'll want to know: \te{కిరాయి}}
\textbf{\te{కిరాయి}} — \emph{kirāyi} — \textquotedblleft{}a fare; hire charge\textquotedblright{}.

\begin{quote}
\te{ఆటో} \te{కిరాయి} \te{ఎంత}?
\end{quote}

\begin{quote}
\emph{āṭō kirāyi enta?}
\end{quote}

\begin{quote}
How much is the auto fare?
\end{quote}

It is the charge for hiring a vehicle, and also the rent of a room or a house.

\begin{cousinweb}
\textbf{\te{కిరాయి}} (\emph{kirāyi}) is a loan from Arabic \emph{kirāʾ}, \textquotedblleft{}hire, rent\textquotedblright{}, which reached Telugu through Persian and Hindustani \emph{kirāya}.
\end{cousinweb}

\begin{grammarlens}[title={What you've built}]
One more word for getting around.
\end{grammarlens}

\subsection*{Your turn}
\begin{itemize}
\raggedright
  \item \emph{Say it:} \emph{kirāyi}
  \item \emph{Say it:} \emph{kirāyi}, once more
  \item \emph{Say it:} \emph{āṭō kirāyi enta?} — How much is the auto fare?
  \item \emph{Recall:} say the Telugu for a sleeper berth, then the Telugu for a railway coach, then say \emph{kirāyi} again
\end{itemize}

\begin{tcolorbox}[breakable,colback=teal!4,colframe=teal!35!black,title={Before you move on}]
What does \te{కిరాయి} mean? (\textquotedblleft{}A fare; hire charge\textquotedblright{}.) Say it once more.
\end{tcolorbox}
